\documentclass[11pt,a4paper]{article}
\usepackage[margin=25mm]{geometry}
\usepackage{amsmath,amssymb,amsthm,mathtools}
\usepackage{graphicx}
\usepackage{xcolor}
\usepackage{microtype}
\usepackage[colorlinks=true,linkcolor=blue!55!black,urlcolor=blue!55!black,citecolor=blue!55!black]{hyperref}

\newtheorem{theorem}{Theorem}
\newcommand{\R}{\mathbb R}
\newcommand{\Pjet}{\mathcal P^2_0(\R^3)}
\newcommand{\Allow}{\operatorname{Allow}}

\title{An implied $2$-jet that is not $1$-implied}
\author{Candidate full counterexample to Remark 2.11 of arXiv:2210.08687}
\date{12 August 2026}

\begin{document}
\maketitle

\begin{center}
\fcolorbox{blue!55!black}{blue!5}{\parbox{0.91\textwidth}{
\textbf{Status: candidate full counterexample; likely valid; human review needed.}
In the truncated $2$-jet ring $\mathcal P^2_0(\mathbb R^3)$, the ideal
$I=\operatorname{span}\{xz,yz,x^2+y^2\}$ implies
$p=x^2+2y^2$, but no single element of $I$ can imply $p$ in the sense of
Remark~2.11.}}
\end{center}

\section*{Source question}
Fefferman and Shaviv define a jet $p$ to be \emph{implied} by an ideal $I$
when, near the allowed directions $\Omega=\Allow(I)$ and at every sufficiently
small scale, $p$ is an arbitrarily small $C^m$ error plus a tame
$C^m$-coefficient combination of finitely many fixed elements of $I$.
Remark~2.11 asks whether the finite list can always be replaced by a single
fixed element of $I$; equivalently, whether every implied jet is
$1$-implied~\cite[p.~12]{FS}.  The exact definition and question are reproduced
below.

\begin{center}
\includegraphics[width=0.98\textwidth]{figures/open_problem_crop.png}
\end{center}

\section*{Result and proof intuition}
All products below are products in the truncated, nonunital jet ring
$\Pjet$: terms of degree greater than $2$ are discarded.

\begin{theorem}\label{thm:main}
Let
\[
 I=\operatorname{span}_{\R}\{xz,yz,x^2+y^2\}\subset\Pjet,
 \qquad p=x^2+2y^2.
\]
Then $I$ is an ideal, $p$ is implied by $I$, and $p$ is not $1$-implied by
$I$.  Consequently the question in Remark~2.11 has a negative answer.
\end{theorem}

\paragraph{Proof intuition.}
The allowed directions of $I$ are precisely the two poles of the unit sphere.
Near those poles, two generators recover $p$ by the elementary identity
\[
 p=\frac{x}{z}(xz)+\frac{2y}{z}(yz).
\]
The apparent freedom of using two generators is essential.  A single
$Q\in I$ has the form $Q=axz+byz+c(x^2+y^2)$.  At a pole, a first derivative
in the transverse $(a,b)$ direction forces its coefficient $S$ to be
arbitrarily small.  A second derivative in the perpendicular transverse
direction then cannot reproduce the nonzero Hessian of $p$.  When
$(a,b)=0$, subtracting the $xx$ and $yy$ second-derivative identities cancels
$S$ and gives the contradiction directly.

\section*{Proof}
Every element of $I$ is homogeneous of degree $2$.  Multiplying it by any
element of the nonunital ring $\Pjet$ produces only terms of degree at least
$3$, hence zero after truncation.  Thus $I$ is an ideal.  Also $p\notin I$:
in a representation of $p$ by the displayed generators, the coefficient of
$x^2+y^2$ would have to be both $1$ (from $x^2$) and $2$ (from $y^2$).

We first determine the allowed directions.  If
$\omega=(\omega_x,\omega_y,\omega_z)\in S^2$ is not a pole, then
$\omega_x^2+\omega_y^2>0$.  By continuity, on a conic neighborhood of
$\omega$ one has
\[
 x^2+y^2\ge c\lvert(x,y,z)\rvert^2
\]
for some $c>0$.  Since $x^2+y^2\in I$, this makes $\omega$ forbidden in the
sense of Definition~2.6.  Conversely, every element of $I$ vanishes on the
punctured $z$-axis.  No finite family from $I$ can therefore satisfy the
required positive lower bound on a cone containing either pole.  Hence
\begin{equation}\label{eq:omega}
 \Omega=\Allow(I)=\{(0,0,1),(0,0,-1)\}.
\end{equation}

We next show that $I$ implies $p$.  Choose a smooth, degree-zero angular
cutoff $\chi$ on $\R^3\setminus\{0\}$ which equals $1$ in conic
neighborhoods of both poles and is supported where
$|z|\ge \tfrac13|(x,y,z)|$.  On the punctured space define, with the evident
smooth zero extension off the support of $\chi$,
\[
 S_1=\chi\frac{x}{z},\qquad S_2=2\chi\frac{y}{z},\qquad F=0,
 \qquad Q_1=xz,\quad Q_2=yz.
\]
The functions $S_1,S_2$ are smooth and homogeneous of degree zero.  Thus,
for every multi-index $|\alpha|\le2$, there is a fixed constant $A$ such that
on $\operatorname{Ann}_4(\rho)$,
\[
 |\partial^\alpha S_j|\le A\rho^{-|\alpha|}\qquad(j=1,2).
\]
After reducing the angular parameter $\delta>0$ if necessary, $\chi=1$
whenever $\operatorname{dist}(X/|X|,\Omega)<\delta$.  There, at every scale,
\[
 F+S_1Q_1+S_2Q_2
 =\frac{x}{z}(xz)+\frac{2y}{z}(yz)=p.
\]
The zero function satisfies all the $\varepsilon$-small bounds required of
$F$.  These fixed $A,Q_1,Q_2$, with any sufficiently small radius $r$, verify
Definition~2.10 for every $\varepsilon>0$.

It remains to rule out $1$-implication.  Suppose to the contrary that a fixed
$Q\in I$ and a constant $A$ witness it.  Write
\begin{equation}\label{eq:q}
 Q=axz+byz+c(x^2+y^2).
\end{equation}
For every $\varepsilon>0$, Definition~2.10 would then supply $\delta,r>0$;
for each $0<\rho\le r$ it would supply $F,S\in C^2(\operatorname{Ann}_4(\rho))$
such that
\begin{equation}\label{eq:identity}
 p=F+SQ
\end{equation}
on the open part of $\operatorname{Ann}_2(\rho)$ whose directions are within
$\delta$ of $\Omega$, and
\begin{equation}\label{eq:fsmall}
 |\partial^\alpha F|\le\varepsilon\rho^{2-|\alpha|}
 \qquad(|\alpha|\le2).
\end{equation}
Fix such a $\rho$ and evaluate derivatives at the pole point
$v=(0,0,\rho)$, an interior point of the equality region.

First suppose $d=(a^2+b^2)^{1/2}>0$.  Let
\[
 e=(a/d,b/d,0),\qquad t=(-b/d,a/d,0).
\]
At $v$ one has
\[
 Q=0,\quad D_eQ=\rho d,\quad D_tQ=0,\quad D_t^2Q=2c.
\]
Since $D_ep(v)=0$, differentiating \eqref{eq:identity} in the $e$ direction
and using \eqref{eq:fsmall} gives
\begin{equation}\label{eq:sbound}
 |S(v)|\rho d=|D_eF(v)|\le\sqrt2\,\varepsilon\rho,
 \qquad |S(v)|\le\frac{\sqrt2\,\varepsilon}{d}.
\end{equation}
Differentiating twice in the $t$ direction eliminates all product-rule terms
except $S(v)D_t^2Q(v)$, so
\[
 D_t^2p(v)=D_t^2F(v)+2cS(v).
\]
Here $D_t^2p=2t_x^2+4t_y^2\ge2$, while the componentwise bounds in
\eqref{eq:fsmall} imply $|D_t^2F(v)|\le2\varepsilon$.  Consequently
\[
 2\le 2\varepsilon+\frac{2\sqrt2|c|}{d}\varepsilon,
\]
which is impossible after choosing $\varepsilon>0$ sufficiently small.

Finally suppose $a=b=0$, including the possibility $Q=0$.  At $v$,
$Q=D_xQ=D_yQ=0$, and differentiating \eqref{eq:identity} twice yields
\[
 2=F_{xx}(v)+2cS(v),\qquad
 4=F_{yy}(v)+2cS(v).
\]
Subtracting gives
$2=F_{yy}(v)-F_{xx}(v)$, whereas \eqref{eq:fsmall} makes the absolute value
of the right-hand side at most $2\varepsilon$.  Taking
$\varepsilon<1$ is the final contradiction.  Thus no $Q\in I$ can witness
$1$-implication, and the theorem follows. \qed

\section*{Scope and novelty note}
The example answers exactly the implication-versus-$1$-implication question
in Remark~2.11; it does not address the analogous variants mentioned there
for Definitions~3.11 and~3.13.  The run's four cheap indexes and targeted
searches through 11 August 2026 for the exact terminology, the paper title,
and follow-up work found the source paper and its sequel but no resolution of
this question.  Novelty is therefore provisional pending expert literature
review.

\begin{thebibliography}{9}
\bibitem{FS}
C.~Fefferman and A.~Shaviv,
\emph{A property of ideals of jets of functions vanishing on a set},
arXiv:2210.08687 (2022).
\end{thebibliography}

\end{document}
